\documentclass[11pt]{article}

\usepackage{amsmath,amssymb,graphicx}
\usepackage[margin=1in]{geometry}
\usepackage{booktabs}

\title{Dirac--K\"ahler Collision Texture and Grade-Exchange Baseline on a Kernel-Induced Cochain Operator Architecture}
\author{Tomislav Rupi\'c \\ Independent Researcher}
\date{March 2026}

\begin{document}

\maketitle

\begin{abstract}
Phase III of the HAOS-IIP program opens only after a strong Phase I--II boundary has been mapped. Frozen pre-geometry did not bind. Weak scalar reinforcement, mesoscopic reinforcement, directional reinforcement, mixed scalar-vector reinforcement, selective interaction gating, and delayed nonlocal memory all remained persistence-null despite bounded nonzero activity. Phase III therefore changes operator class rather than feedback law. Stage 23 replaces the scalar / projected-transverse propagation sector with a discrete Dirac--K\"ahler sector and asks a narrower question first: does graded first-order propagation change collision phenomenology in ways relevant to persistence before any binding claim is made? Stage 23.1 runs a 10-case collision geometry scan, then corrects its initial read by reclassifying encounters from the combined field rather than from independently propagated packet proxies. Under that stricter read the Dirac--K\"ahler sector produces no candidate bounded collision regimes. A matched scalar control matrix shows that the strongest clustered persistence signal is not uniquely Dirac--K\"ahler. Clustered-pair refinement to $n=24$ also removes the apparent out-of-phase anomaly as a persistence candidate. Stage 23.2 then asks a different question: whether the Dirac--K\"ahler sector supports a structured family of grade-exchange and collision-texture behaviors even when persistence remains weak. The 12-run follow-up across corridor, glancing, and clustered geometries at $n=12$ and $n=24$ shows that pass-through with grade exchange is a real Dirac--K\"ahler collision class, that strong grade-exchange amplitudes up to $2.303\times10^{-1}$ appear in several $n=12$ geometries, and that one family, the in-phase tight clustered pair, remains refinement-stable as a mixed clustered encounter. The bounded conclusion is therefore not that Dirac--K\"ahler propagation binds. It is that Dirac--K\"ahler propagation opens a new collision-texture layer through structured grade exchange, while a persistence advantage over the matched scalar control is not yet established.
\end{abstract}

\section{Why Phase III Exists}

Phase III is justified by exclusion, not enthusiasm. Phase I closed with a hard statement: frozen pre-geometry can organize, but it does not bind. Phase II then tested the largest nearby family of weak mechanism changes that could plausibly have created persistence without changing the packet ontology itself. Local scalar substrate response, mesoscopic basin reinforcement, transport-corridor reinforcement, directional flux reinforcement, mixed scalar-vector reinforcement, selective registration gating, and delayed nonlocal memory all entered the architecture cleanly and remained persistence-null.

This matters because it removes the simplest nearby explanations. The next lawful move is therefore not another weak reinforcement variant. It is an operator-class change. The Dirac--K\"ahler sector is the minimal strong candidate because it changes the state algebra itself. In place of scalar or projected-transverse propagation it introduces graded cochain states and first-order transport via
\[
D_{\mathrm{DK}} = d + \delta.
\]
The Phase III opener does not assume fermionic binding, exclusion, or matter formation. It asks whether graded propagation changes the collision atlas at all.

\section{Stage 23.1: Corrected Dirac--K\"ahler Collision Scan}

Stage 23.1 runs a minimal 10-case geometry-of-encounter scan at $n=12$ on a regular periodic lattice. Five geometries are used with in-phase and out-of-phase initial relations:
\begin{itemize}
\item head-on symmetric pair,
\item counter-propagating corridor pair,
\item offset glancing collision,
\item tight clustered pair,
\item symmetric triad collision.
\end{itemize}
The packet state is minimal graded support on $0$- and $1$-forms, with no adaptive kernel, no memory law, and no extra mediator field.

The important methodological point is that the first raw read of Stage 23.1 was too optimistic. Collision labels inferred from independently propagated packet proxies are not strict enough in a linear sector. The scan was therefore rerun with encounter geometry extracted from the combined field through peak tracking on the total graded intensity. Only the corrected read is retained here.

Under that corrected read the 10-case Dirac--K\"ahler scan returns:
\begin{itemize}
\item $0/10$ candidate bounded collision regimes,
\item $2/10$ weak persistence gains,
\item $8/10$ runs with no persistence gain,
\item $9/10$ runs labeled ``pass-through with grade exchange,''
\item $1/10$ run labeled ``unresolved / mixed.''
\end{itemize}

The strongest clustered-pair cases remain informative but not decisive. The in-phase clustered case produces a long close-packed encounter with composite lifetime $0.7517$, binding persistence $1.0$, and grade-transfer amplitude $1.303\times10^{-1}$, yet does not satisfy the bounded-collision gate. The out-of-phase clustered case produces the largest Stage 23.1 grade-transfer amplitude, $2.303\times10^{-1}$, but only weak persistence gain and no bounded post-collision regime.

The resulting Stage 23.1 statement is therefore narrow:
\begin{quote}
Dirac--K\"ahler propagation changes collision texture materially, but Stage 23.1 alone does not establish a persistence-level regime change.
\end{quote}

\section{Matched Scalar Control}

The matched scalar control is essential because the clustered Dirac--K\"ahler anomaly could otherwise be overstated. The same 10 frozen geometry-phase cases were rerun with a scalar bosonic propagator on the same periodic $2$D complex. The control returns:
\begin{itemize}
\item $1/10$ candidate bounded collision regime,
\item $9/10$ runs with no persistence gain,
\item one in-phase tight clustered pair labeled ``metastable composite.''
\end{itemize}

The control clustered in-phase case reaches composite lifetime $0.6058$ and binding persistence $1.0$, which is enough to block any claim that Dirac--K\"ahler propagation has already opened a unique persistence channel at this stage. Put differently, the strongest clustered persistence signature in the opener scan is not uniquely Dirac--K\"ahler.

This comparison is the reason the Phase III opener must be stated in terms of collision phenomenology rather than early binding.

\section{Clustered-Pair Refinement}

To test whether the Dirac--K\"ahler clustered family strengthens under refinement, the clustered pair alone was promoted to $n=24$. The refinement result is again informative but bounded.
\begin{itemize}
\item the in-phase clustered case remains an ``unresolved / mixed'' close-packed encounter,
\item the persistence label remains only ``weak persistence gain,''
\item the out-of-phase clustered case falls back to ``pass-through dispersive,''
\item no candidate bounded collision regime appears.
\end{itemize}

The in-phase clustered case remains long-lived as an encounter, with composite lifetime $0.7489$ and binding persistence $1.0$, but its grade-transfer amplitude falls to $4.140\times10^{-2}$. The out-of-phase clustered case loses even the earlier weak anomaly. This is strong evidence that Stage 23.1 did not discover a robust persistence family.

\section{Stage 23.2: Grade-Exchange and Collision-Texture Scan}

Stage 23.2 follows the signal that survived comparison. It does not ask whether the Dirac--K\"ahler sector binds. It asks whether the sector supports a structured family of grade-exchange and collision-texture behaviors.

The scan is kept narrow:
\begin{itemize}
\item corridor pair,
\item glancing collision,
\item tight clustered pair,
\item in-phase and out-of-phase initial relations,
\item both $n=12$ and $n=24$.
\end{itemize}
This gives a 12-run follow-up focused on the geometries that were most informative in Stage 23.1.

The Stage 23.2 summary is:
\begin{itemize}
\item $5$ runs labeled ``pass-through with grade exchange,''
\item $5$ runs labeled ``pass-through dispersive,''
\item $2$ runs labeled ``unresolved / mixed,''
\item $3$ runs with weak persistence gain,
\item $9$ runs with no persistence gain,
\item one refinement-stable texture family.
\end{itemize}

The grade-exchange signatures separate into five classes:
\begin{itemize}
\item clustered mixed encounter ($2$ runs),
\item strong pass-through exchange ($4$ runs),
\item weak pass-through exchange ($1$ run),
\item dispersive exchange tail ($1$ run),
\item texture null ($4$ runs).
\end{itemize}
The strongest exchange amplitude in the whole follow-up remains the Stage 23.1-style clustered out-of-phase value $2.303\times10^{-1}$ at $n=12$, but that family is not stable under refinement. The only family that remains stable under $n=12\rightarrow24$ is the in-phase tight clustered pair, which stays in the signature class ``clustered mixed encounter.''

\section{What Stage 23.2 Adds}

Stage 23.2 does not rescue a persistence claim. What it adds is structure. Three points matter.

\subsection{Grade Exchange Is Not Noise}

Pass-through with grade exchange is not a one-off artifact. It appears repeatedly across corridor, glancing, and clustered geometries at $n=12$, with amplitudes from roughly $1.48\times10^{-1}$ up to $2.30\times10^{-1}$. This already distinguishes the Dirac--K\"ahler sector from a plain scalar dispersive story.

\subsection{Phase Sensitivity Exists But Is Not Uniform}

The out-of-phase clustered case is the strongest short-scale exchange event, while the in-phase clustered case is the only family that remains stable under refinement as a mixed encounter. Corridor and glancing cases carry strong exchange at $n=12$ but relax toward dispersive labels at $n=24$. So the phase relation matters, but not as a simple monotonic gate.

\subsection{One Family Survives Refinement as Texture, Not Binding}

The in-phase tight clustered pair is the first clear refinement-stable Phase III family. Its stability is not persistence-level boundedness. It is stability of collision texture: a repeated close-packed mixed encounter with weak persistence gain and nonzero grade transfer at both resolutions.

\section{Phase III Opener Boundary}

The data now support a cleaner and stronger opener statement than Stage 23.1 alone allowed.

\begin{table}[h]
\centering
\small
\begin{tabular}{@{}p{2.2in}p{1.4in}p{1.0in}p{1.4in}@{}}
\toprule
Probe & Strongest signal & Persistence status & Bounded conclusion \\
\midrule
23.1 corrected DK base scan & widespread pass-through grade exchange & no bounded regime & DK changes collision texture \\
Matched scalar control & clustered in-phase metastable composite & stronger than DK on persistence at $n=12$ & DK persistence advantage not established \\
23.1 clustered refinement & in-phase mixed clustered encounter survives & weak only & out-of-phase anomaly washes out \\
23.2 texture scan & one stable clustered in-phase texture family & weak only & DK supports structured exchange families \\
\bottomrule
\end{tabular}
\caption{Phase III opener logic. The Dirac--K\"ahler sector earns a collision-phenomenology claim, not yet a persistence claim.}
\end{table}

The strongest compact statement is therefore:
\begin{quote}
Dirac--K\"ahler propagation opens a new collision-texture layer through structured grade exchange, while a persistence advantage over the matched scalar control is not yet established.
\end{quote}

This is a real positive signal, but it is narrower than binding. The operator sector now appears to matter for encounter structure, grade transfer, and phase-sensitive collision classes. It does not yet matter in a decisive way for persistence.

\section{Conclusion}

Phase III opens successfully, but not on the axis first hoped for. The Dirac--K\"ahler sector does not immediately generate bound composites or a clean persistence gain over the matched scalar control. What it does generate is a new phenomenology of collisions: widespread pass-through with grade exchange, geometry-dependent exchange amplitudes, and one refinement-stable clustered mixed-encounter family.

That is enough to justify continuation of Phase III. The next lawful move is not a generic search for binding. It is a focused search for Dirac--K\"ahler-specific collision structure, exchange channels, and refinement-stable graded encounter families.

The resulting interpretation boundary remains strict. Phase III has opened a new operator-sector atlas. It has not yet opened a new persistence law.

\end{document}
